\documentclass[a4paper, 12pt]{article}

\usepackage[portuges]{babel}
\usepackage[utf8]{inputenc}
\usepackage{amsmath}
\usepackage{indentfirst}
\usepackage{graphicx}
\usepackage[colorinlistoftodos]{todonotes}
\usepackage{verbatim}
\usepackage{textcomp}
\usepackage{gensymb}

\begin{document}

\begin{titlepage}
\begin{center}
\textbf{\LARGE Universidade de Brasília}\\[0.5cm] 
\textbf{\large Departamento de Engenharia Elétrica}\\[0.2cm]
\vspace{20pt}
\includegraphics{Logo_UnB.png}\\[1cm]

\par
\vspace{32pt}
\textbf{\LARGE Projeto Final do Laboratório de Circuitos Elétricos 2}\\
\vspace{30pt}
\textbf {\Large Autores:}\\[0.2cm]
\Large {André Bauer		11/0055055}\\[0.1cm]
\Large {Felipe Brandão	12/0044919}\\[0.1cm]
\Large {Igor Miranda	11/0013255}\\[0.1cm]

\end{center}

\par
\vfill
\begin{center}
{{\normalsize Brasília}\\
{\normalsize \today}}
\end{center}
\end{titlepage}

%Sumário
\newpage
\tableofcontents
\thispagestyle{empty}
%End Sumário

\newpage
\section{Resumo}



\section{Introdução}

\subsection{Introdução Teórica}


\begin{gather*}
p(t) = v(t)\cdot i(t)\\
\end{gather*}

\subsection{Contextualização}


\begin{figure}[h]
	\centering
    \includegraphics{Circuito.png}
    \caption{Circuito a ser projetado}
\end{figure}

\section{Cálculos Teóricos}

\subsection{Potência Instantânea e Potência Média}


\begin{gather*}
\textbf{V} = R\cdot\textbf{I}\\
\end{gather*}


\begin{multline*}
p(t) = \frac{1}{2}\cdot499,91m\cdot227,23\mu\cdot\cos(1,08\degree - 1,08\degree) +\\
\frac{1}{2}\cdot499,91m\cdot227,23\mu\cdot\cos(2\cdot3,14k\cdot t + 1,08\degree + 1,08\degree)
\end{multline*}

$$p(t) \approx 56,80\mu + 56,80\mu\cdot\cos(6,28k\cdot t + 2,16\degree)\;W$$



\subsection{Potência Complexa}


\begin{gather*}
\textbf{V}_{RMS} = \frac{\textbf{V}}{\sqrt{2}}\\
\end{gather*}

\section{Simulação Computacional}


\begin{figure}[ht]
	\centering
    \includegraphics[width=\textwidth]{Simula__o.PNG}
    \caption{}
\end{figure}

\section{Experimento}


\begin{enumerate}
	\item
\end{enumerate}



\section{Conclusão}



\section{Referência}

\noindent

\end{document}